\documentclass[10pt,a4paper]{amsart}
\usepackage[margin=19mm]{geometry}
\usepackage{amsmath,amssymb,mathtools}
\usepackage[scr=rsfs,cal=euler]{mathalfa}
\usepackage{xfrac}
\usepackage[unicode,hidelinks,bookmarks=false]{hyperref}
\newcommand{\ie}{i.e.\ }
\newcommand{\cf}{cf.\ }
\newcommand{\resp}{respectively\ }
\newcommand{\Subsection}{\subsection}
\providecommand{\nobreakdash}{\nobreak\hbox{-}}
\setlength{\emergencystretch}{2em}
\multlinegap0pt
\usepackage{bm}
\usepackage{bbm}
\usepackage{dsfont}
\usepackage{xspace}
\usepackage{cancel}
\usepackage{mathtools}
\usepackage{amsthm}
\makeatletter
\@ifundefined{lemma}{\newtheorem{lemma}{Lemma}}{}
\@ifundefined{theorem}{\newtheorem{theorem}{Theorem}}{}
\makeatother
\title[Generating the Spanning Trees of Series-Parallel Graphs up t]{Generating the Spanning Trees of Series-Parallel Graphs up to Graph Automorphism}
\author{Mithra Karamchedu, Lucas Bang}
\date{Source: arXiv:2508.13480v1 (CC BY 4.0)}
\begin{document}
\maketitle

\noindent\emph{Source and licence.} Generating the Spanning Trees of Series-Parallel Graphs up to Graph Automorphism. Version arXiv:2508.13480v1 (CC BY 4.0), distributed under the Creative Commons Attribution 4.0 International licence, as recorded in the arXiv licence metadata for this exact version and shown on the abstract page https://arxiv.org/abs/2508.13480v1. Full rights notice: licenses/arxiv-2508.13480-CC-BY-4.0.txt.\par\smallskip

\noindent\textit{Editorial scope.} Counted unit: the Theorem whose source label is \texttt{runtime} in the article (source0.tex lines 946--949, source label \texttt{runtime}), delivered together with the article's own complete original proof of it (source0.tex lines 950--999, one \texttt{proof} environment of 50 lines). No mathematics is written, repaired or re-derived: statement and proof are the article's own text line for line. Required context retained uncounted: seriallemma1 (lines 895--900); seriallemma2 (lines 912--917); parallellemma (lines 926--931). Every definition, display and label of those lines is retained unchanged. Omitted from this package and not redistributed: 1--894, 901--911, 918--925, 932--945, 1000--1214. Nothing inside a retained excerpt is omitted or altered: every retained body is delivered exactly as the article's lines are written, so no omission receipt of any kind accompanies this package. Citations: the retained excerpts carry no citation command at all, so no bibliography entry of the article is transcribed and no reference is redistributed. Reference adaptation: none was needed and none was made; every reference inside the delivered unit resolves inside it, so no unresolved reference or citation is produced. Printed slips kept literal: no printed word, symbol, hypothesis or number of the counted statement or of its proof was altered. Layout and class: the article's own class is \texttt{siamart171218} with packages including geometry, tikz, circuitikz, amsmath, amssymb, ifthen, lipsum, amsfonts, graphicx, epstopdf, algorithmic, amsopn; the wrapper is the lane's pinned \texttt{amsart} editorial wrapper at 10pt on A4, which restates the theorem-like environments the retained text needs and adds the article's own macro definitions that the retained text needs (none), with extra wrapper packages (bm, bbm, dsfont, xspace, cancel, mathtools). The delivered numbering is the unit's own. Assets: no retained span contains a figure environment, a graphics-inclusion command, a PDF-page inclusion, a table or a TikZ picture; no image file is copied into this package, the package declares no asset dependency and no image file is redistributed with it. Licence: version arXiv:2508.13480v1 of this article is distributed under the Creative Commons Attribution 4.0 International licence, as recorded in the arXiv licence metadata for this exact version and shown on the abstract page https://arxiv.org/abs/2508.13480v1. A full rights notice is delivered at licenses/arxiv-2508.13480-CC-BY-4.0.txt. Characters: every retained excerpt is pure ASCII, so the delivered engine, XeLaTeX with its default Latin Modern text font, reports no missing character.
\begin{lemma}
    \label{seriallemma1}
    Let $G$ be an oriented serial superedge, such that each $G_i$ is non-serial. Then,
    \[\sum_{i = 1}^k n_i\left\vert \frac{\emph{\textsf{ST}}(G_i)}{\emph{Aut}_\emph{or}(G_i)}\right\vert = O\left(n\left\vert \frac{\emph{\textsf{ST}}(G)}{\emph{Aut}_\emph{or}(G)}\right\vert \right),\]
    where $G_i$ has $n_i$ vertices, for $1 \leq i \leq k$.
\end{lemma}

\begin{lemma}
    \label{seriallemma2}
    Let $G$ be an oriented serial superedge, such that each $G_i$ is non-serial. Then,
    \[\sum_{i = 1}^k n_i\left\vert \frac{\emph{\textsf{NT}}(G_i)}{\emph{Aut}_\emph{or}(G_i)}\right\vert = O\left(n\left\vert \frac{\emph{\textsf{NT}}(G)}{\emph{Aut}_\emph{or}(G)}\right\vert \right),\]
    where $G_i$ has $n_i$ vertices, for $1 \leq i \leq k$.
\end{lemma}

\begin{lemma}
    \label{parallellemma}
    Let $G$ be an oriented parallel superedge, such that each $G_i$ is non-parallel. Then,
    \[\sum_{i = 1}^\ell n_i\left[\left\vert \frac{\emph{\textsf{ST}}(R_i)}{\emph{Aut}_\emph{or}(G_i)}\right\vert + \left\vert \frac{\emph{\textsf{NT}}(R_i)}{\emph{Aut}_\emph{or}(G_i)}\right\vert\right] = O\left(n\left\vert \frac{\emph{\textsf{ST}}(G)}{\emph{Aut}_\emph{or}(G)}\right\vert\right),\]
    where the class representative $R_i$ has $n_i$ vertices, for $1 \leq i \leq k$. 
\end{lemma}

\begin{theorem}
    \label{runtime}
    If $G$ is a series-parallel graph with $n$ vertices and $m$ edges, then \[B(G) = \Theta\left(n\left[\left\vert \frac{\emph{\textsf{ST}}(G)}{\emph{Aut}_\emph{or}(G)}\right\vert + \left\vert \frac{\emph{\textsf{NT}}(G)}{\emph{Aut}_\emph{or}(G)}\right\vert\right]\right) \qquad\text{and}\qquad T(G) = \Theta\left(n\left\vert \frac{\emph{\textsf{ST}}(G)}{\emph{Aut}_\emph{or}(G)}\right\vert\right).\]
\end{theorem}

\begin{proof}
    To prove this result, we will apply strong induction over the depth $d$ of our series-parallel decomposition tree.
    \vspace{5pt}
    \begin{enumerate}
        \item \underline{Base Case}:\\[6pt]
        Let us consider the case where $d = 0$. Then, $G$ is a single edge with endpoints $s$ and $t$. In this case, \texttt{serial\_both\_trees} and \texttt{parallel\_both\_trees} return the single nonequivalent spanning tree of $G$ ($G$ itself) and near tree ($\overline{G}$) in constant time $\Theta(1)$. Since $G$ is a single edge, $n = 2$ and $\vert \textsf{ST}(G)/\text{Aut}_\text{or}(G)\vert = \vert \textsf{NT}(G)/\text{Aut}_\text{or}(G)\vert = 1$, so $\Theta\big(n\cdot [\vert \textsf{ST}(G)/\text{Aut}_\text{or}(G)\vert + \vert \textsf{NT}(G)/\text{Aut}_\text{or}(G)\vert] \big)$ also equals $\Theta(1)$. Hence,
        \[B(G) = \Theta\left(n\left[\left\vert \frac{\text{\textsf{ST}}(G)}{\text{Aut}_\text{or}(G)}\right\vert + \left\vert \frac{\text{\textsf{NT}}(G)}{\text{Aut}_\text{or}(G)}\right\vert\right]\right).\]
        By similar reasoning, $T(G) = \Theta\big(n\cdot \vert \textsf{ST}(G) / \text{Aut}_\text{or}(G)\vert\big)$. Thus, the claim holds when $d = 0$.\\
        \item \underline{Inductive Hypothesis}:\\[6pt]
        Let us assume that the claim is true for all $d \leq h$, for some $h \in \mathbb{N}$. That is, for all series-parallel graphs $G$ with depth  at most $h$, we have that $B(G) = \Theta\big(n\cdot [\vert \textsf{ST}(G)/\text{Aut}_\text{or}(G)\vert + \vert \textsf{NT}(G)/\text{Aut}_\text{or}(G)\vert]\big)$ and $T(G) = \Theta\big(n\cdot \vert \textsf{ST}(G)/\text{Aut}_\text{or}(G)\vert\big)$. We will show that the claim is then true for depth $d = h + 1$.\\

        \item \underline{Induction Step}:\\[6pt]
        Let $G$ be any series-parallel graph with depth $d = h + 1$. Then, the child subtrees of $G$ (which are the decomposition trees of each $G_i$) each have depth at most $h$. If $G$ is serial, then we have that
        \begin{align*}
            B_S(G) &= \sum_{i = 1}^k B(G_i) + \Theta\left(n\left[\left\vert \frac{\textsf{ST}(G)}{\text{Aut}_\text{or}(G)}\right\vert + \left\vert \frac{\textsf{NT}(G)}{\text{Aut}_\text{or}(G)}\right\vert\right]\right)\\[6pt]
            &= \Theta\left(n\left[\left\vert \frac{\textsf{ST}(G)}{\text{Aut}_\text{or}(G)}\right\vert + \left\vert \frac{\textsf{NT}(G)}{\text{Aut}_\text{or}(G)}\right\vert\right]\right) + \sum_{i = 1}^k \Theta\left(n_i\left[\left\vert \frac{\textsf{ST}(G_i)}{\text{Aut}_\text{or}(G_i)}\right\vert + \left\vert \frac{\textsf{NT}(G_i)}{\text{Aut}_\text{or}(G_i)}\right\vert\right]\right)\\[6pt]
            &= \Theta\left(n\left[\left\vert \frac{\textsf{ST}(G)}{\text{Aut}_\text{or}(G)}\right\vert + \left\vert \frac{\textsf{NT}(G)}{\text{Aut}_\text{or}(G)}\right\vert\right] + \sum_{i = 1}^k n_i\left[\left\vert \frac{\textsf{ST}(G_i)}{\text{Aut}_\text{or}(G_i)}\right\vert + \left\vert \frac{\textsf{NT}(G_i)}{\text{Aut}_\text{or}(G_i)}\right\vert\right]\right),
        \end{align*}
        where each $G_i$ has $n_i$ vertices. By Lemmas \ref{seriallemma1} and \ref{seriallemma2}, $\sum_{i = 1}^k n_i\cdot [\vert \textsf{ST}(G_i)/\text{Aut}_\text{or}(G_i)\vert] = O\big(n\cdot\vert \textsf{ST}(G)/\text{Aut}_\text{or}(G)\vert\big)$ and $\sum_{i = 1}^k n_i\cdot [\vert \textsf{NT}(G_i)/\text{Aut}_\text{or}(G_i)\vert] = O\big(n\cdot\vert \textsf{NT}(G)/\text{Aut}_\text{or}(G)\vert\big)$. Therefore, the summation term is at most the same order as the leading term, so we deduce that
        \[B_S(G) = \Theta\left(n\left\vert \frac{\textsf{ST}(G_i)}{\text{Aut}_\text{or}(G_i)}\right\vert\right).\]
        If $G$ is parallel, then we proceed similarly:
        \begin{align*}
            B_P(G) &= \sum_{i = 1}^\ell B(R_i) + \Theta\left(n\left[\left\vert \frac{\textsf{ST}(G)}{\text{Aut}_\text{or}(G)}\right\vert + \left\vert \frac{\textsf{NT}(G)}{\text{Aut}_\text{or}(G)}\right\vert\right]\right)\\[6pt]
            &= \Theta\left(n\left[\left\vert \frac{\textsf{ST}(G)}{\text{Aut}_\text{or}(G)}\right\vert + \left\vert \frac{\textsf{NT}(G)}{\text{Aut}_\text{or}(G)}\right\vert\right]\right) + \sum_{i = 1}^\ell \Theta\left(n_i\left[\left\vert \frac{\textsf{ST}(R_i)}{\text{Aut}_\text{or}(R_i)}\right\vert + \left\vert \frac{\textsf{NT}(R_i)}{\text{Aut}_\text{or}(R_i)}\right\vert\right]\right)\\[6pt]
            &= \Theta\left(n\left[\left\vert \frac{\textsf{ST}(G)}{\text{Aut}_\text{or}(G)}\right\vert + \left\vert \frac{\textsf{NT}(G)}{\text{Aut}_\text{or}(G)}\right\vert\right] + \sum_{i = 1}^\ell n_i\left[\left\vert \frac{\textsf{ST}(R_i)}{\text{Aut}_\text{or}(R_i)}\right\vert + \left\vert \frac{\textsf{NT}(R_i)}{\text{Aut}_\text{or}(R_i)}\right\vert\right]\right),
        \end{align*}
        By Lemma \ref{parallellemma}, the summation term has order at most $O\big(n\cdot\vert \textsf{ST}(G)/\text{Aut}_\text{or}(G)\vert\big)$, so
        \[B_P(G) =  \Theta\left(n\left[\left\vert \frac{\textsf{ST}(G)}{\text{Aut}_\text{or}(G)}\right\vert + \left\vert \frac{\textsf{NT}(G)}{\text{Aut}_\text{or}(G)}\right\vert\right] \right).\] Combining our asymptotic expressions for $B_S(G)$ and $B_P(G)$ gives us that
        \[B(G) = \Theta\left(n\left[\left\vert \frac{\textsf{ST}(G)}{\text{Aut}_\text{or}(G)}\right\vert + \left\vert \frac{\textsf{NT}(G)}{\text{Aut}_\text{or}(G)}\right\vert\right] \right),\] as desired.

        Let us now consider $T(G)$. If $G$ is serial, then we have that
        \begin{align*}
            T_S(G) = \sum_{i = 1}^k T(G_i) + \Theta\left(n\left\vert \frac{\textsf{ST}(G)}{\text{Aut}_\text{or}(G)}\right\vert\right)  &= \Theta\left(n\left\vert \frac{\textsf{ST}(G)}{\text{Aut}_\text{or}(G)}\right\vert\right) + \sum_{i = 1}^\ell \Theta\left(n_i\left\vert \frac{\textsf{ST}(R_i)}{\text{Aut}_\text{or}(R_i)}\right\vert\right)\\[6pt]
             &= \Theta\left(n\left\vert \frac{\textsf{ST}(G)}{\text{Aut}_\text{or}(G)}\right\vert + \sum_{i = 1}^\ell n_i\left\vert \frac{\textsf{ST}(R_i)}{\text{Aut}_\text{or}(R_i)}\right\vert\right).
        \end{align*}
        By Lemma \ref{parallellemma}, the summation term once again has order at most $O\big(n\cdot\vert\textsf{ST}(G)/\text{Aut}_\text{or}(G)\vert\big)$, so as we might expect,
        \[T_S(G) = \Theta\left(n\left\vert \frac{\textsf{ST}(G)}{\text{Aut}_\text{or}(G)}\right\vert\right).\]
        Finally, if $G$ is parallel,
        then
         \begin{align*}
            T_P(G) &= \sum_{i = 1}^\ell B_S(R_i) + \Theta\left(n\left\vert \frac{\textsf{ST}(G)}{\text{Aut}_\text{or}(G)}\right\vert\right).\\[6pt]
             &= \Theta\left(n\left\vert \frac{\textsf{ST}(G)}{\text{Aut}_\text{or}(G)}\right\vert\right) + \sum_{i = 1}^\ell \Theta\left(n_i\left[\left\vert \frac{\textsf{ST}(R_i)}{\text{Aut}_\text{or} (R_i)}\right\vert + \left\vert \frac{\textsf{NT}(R_i)}{\text{Aut}_\text{or} (R_i)}\right\vert\right]\right)\\[6pt]
             &= \Theta\left(n\left\vert \frac{\textsf{ST}(G)}{\text{Aut}_\text{or}(G)}\right\vert + \sum_{i = 1}^\ell n_i\left[\left\vert \frac{\textsf{ST}(R_i)}{\text{Aut}_\text{or} (R_i)}\right\vert + \left\vert \frac{\textsf{NT}(R_i)}{\text{Aut}_\text{or} (R_i)}\right\vert\right]\right).
        \end{align*}
        By Lemma \ref{parallellemma}, the summation term has order at most $O\big(n\cdot \vert \textsf{ST}(G)/\text{Aut}_\text{or}(G)\vert\big)$, so $T_P(G) = \Theta\big(n\cdot \vert \textsf{ST}(G)/\text{Aut}_\text{or}(G)\vert\big)$. Using our asymptotic expressions for $T_S(G)$ and $T_P(G)$, we conclude that
        \[T(G) = \Theta\left(n\left\vert \frac{\textsf{ST}(G)}{\text{Aut}_\text{or}(G)}\right\vert\right),\]
        which is the desired result.
    \end{enumerate}
    By induction, we conclude that \[B(G) = \Theta\left(n\left[\left\vert \frac{\text{\textsf{ST}}(G)}{\text{Aut}_\text{or}(G)}\right\vert + \left\vert \frac{\text{\textsf{NT}}(G)}{\text{Aut}_\text{or}(G)}\right\vert\right]\right) \qquad\text{and}\qquad T(G) = \Theta\left(n\left\vert \frac{\textsf{ST}(G)}{\text{Aut}_\text{or}(G)}\right\vert\right).\]
\end{proof}

\end{document}
